\section{Cronograma de actividades}

A continuación se presenta el cronograma de actividades llevadas a cabo durante el desarrollo del proyecto, 
elaborado en Miro como estimación inicial del proceso. Si bien surgieron inconvenientes que requirieron 
ajustes y cambios de enfoque en algunos aspectos, el cronograma se mantuvo relativamente estable a lo 
largo del desarrollo.

\vspace{0.25cm}

Como herramienta de trabajo se utilizó el modelo USM, donde los requisitos principales se organizaron verticalmente 
y los \textit{features} asociados a cada uno se distribuyeron de forma horizontal, permitiendo definir el backlog 
de cada versión. Ocasionalmente se incorporaron \textit{features} no planificados, ya sea por subestimar el alcance 
de algunos requisitos o por la aparición de nuevas ideas consideradas relevantes para el producto. En caso 
de no completar todos los \textit{features} de una versión, estos se trasladaban a la siguiente con mayor prioridad.

\vspace{0.25cm}

Se planificaron 7 \textit{releases}, desde la \textbf{v0.1.0} hasta la \textbf{v0.7.0}, donde cada una incorporaba un conjunto de 
funcionalidades de complejidad creciente. A continuación se detalla un resumen de lo desarrollado en 
cada \textit{release}.

\subsection{v0.1.0: 28 de octubre, 2025}

Primera versión del proyecto, orientada a establecer las bases técnicas del sistema. Se implementaron 
las funcionalidades esenciales del juego como el inventario, el sistema de stamina, el daño a enemigos 
y el spawn de los mismos. Estos componentes se deprecaron rapidamente a medida que se avanzaba en el desarrollo
y requirieron ser refactorizadas en versiones posteriores, pero sirvieron como punto de partida para validar 
conceptos y establecer la arquitectura inicial. En paralelo, se desarrollaron pruebas de concepto para el modo multijugador 
y la edición de terreno, junto con una interfaz inicial para el editor de mundos y la homepage. 

\vspace{0.25cm}

También se incorporó la autenticación de usuarios mediante la creación e inicio de sesión de cuentas. 
En esta etapa, los componentes del sistema aún no se encontraban integrados entre sí, contando únicamente 
con el cliente de juego, el editor de mundos, el core-service y el Shared Package.

\subsection{v0.2.0: 17 de diciembre, 2025}

Esta versión consolidó la integración entre el editor de mundos y el cliente de juego. Se implementó 
la publicación y carga de mundos, permitiendo que los jugadores pudieran unirse a mundos desde la 
homepage. Se incorporaron funcionalidades de juego como la visualización de bolsas de botín, la 
recolección de ítems, los drops de enemigos y el creador de personajes. Sin embargo, el despliegue 
de mundos aún requería un proceso manual para levantar el servidor dedicado. A nivel de red, se 
utilizaba NetCode como \textit{framework} de multijugador, cuya implementación comenzaba a escalar pero 
presentaba problemas de mantenibilidad. El editor de mundos también crecía en complejidad, aunque 
el sistema de colocación de objetos estaba limitado a una grilla cuadrada con objetos hardcodeados.

\subsection{v0.3.0: 23 de enero, 2026}

Versión centrada en la expansión del contenido jugable y las herramientas de creación. Se incorporó 
el sistema de \textit{Quests}, incluyendo la aceptación, seguimiento y completado de misiones mediante 
la derrota de enemigos e interacción con NPCs. Se agregaron armas, consumibles, tiendas y el sistema 
de diálogos con NPCs. En el editor se habilitó la configuración de spawners de enemigos, jugadores y 
NPCs, tablas de botín y precios de ítems en oro y gemas. A nivel técnico, se migró el servidor dedicado 
de NetCode a Mirror y se refactorizó la arquitectura del editor de mundos, reemplazando el sistema de 
grillas por un sistema de colocación libre de objetos, mejorando significativamente la flexibilidad 
y mantenibilidad. También se creó la landing page del proyecto, con información general,
un blog del desarrollo y links a las redes sociales.

\subsection{v0.4.0: 06 de marzo, 2026}

Durante esta versión no se logró completar el backlog planificado debido a compromisos académicos 
externos que el equipo debió atender, principalmente la preparación de exámenes finales de materias pendientes.
Lo único que se logró resolver fue arreglar bugs en el editor de mundos, y completar refactorizaciones en el servidor
de juego.

\subsection{v0.5.0: 06 de abril, 2026}

Versión dedicada principalmente a la estabilización y refactorización de los sistemas existentes adaptandolos
a la nueva implementación del servidor.
Se reestructuraron los módulos de salud, ataque, dash, stamina, inventario, botín, spawners y diálogos. 
Se integró Nomad para el despliegue dinámico de servidores y Datadog para la observabilidad del sistema. 
Se incorporaron portales para el viaje entre zonas, la tienda de gemas con método de pago y la 
inteligencia artificial de enemigos con comportamiento de persecución y ataque. 

\vspace{0.25cm}

Por primera vez, los creadores pudieron diseñar mundos con múltiples zonas ejecutándose en servidores independientes, 
habilitando mundos de escala ilimitada. También se implementó la lógica de bots para pruebas de carga 
masivas, complementada con la integración de Datadog para la recolección de métricas durante dichas pruebas. 

\subsection{v0.6.0: 10 de mayo, 2026}

Versión avanzada orientada a completar las funcionalidades del producto mínimo viable. Se implementó el 
sistema de cosméticos con inventario global, la persistencia del personaje entre sesiones, el chat 
flotante en multijugador y los nametags sobre los personajes. Se incorporaron cofres, rampas, cambio 
de material de suelo y cielo, música ambiental y efectos de sonido. 

A nivel económico, se completó el sistema de suscripciones, transferencias a creadores y compra de 
cosméticos con gemas.

\vspace{0.25cm}

Esta versión marcó un hito al realizarse la primera demo con personas externas al equipo, cuyo feedback 
resultó valioso para definir las funcionalidades prioritarias del MVP.
También marcó el punto en que el proyecto comenzó a percibirse como un producto consolidado, lo que impactó positivamente 
en la motivación del equipo para llevar el MVP a su estado final.

\subsection{v0.7.0: 05 de junio, 2026}

Versión de cierre del proyecto, orientada a la corrección de bugs, ajustes finales y mejoras de 
calidad de vida (\textit{Quality of Life}). Se integró pruebas unitarias, de integración y de aceptación
en el core service con el fin de aumentar la cobertura de pruebas y mejorar la estabilidad del sistema. 

Se incorporó el Back Office para que los administradores puedan gestionar mundos, 
cuentas de usuarios y contenido de la plataforma. La landing page fue actualizada para incluir 
la descarga del cliente de juego y el editor de mundos.

\vspace{0.25cm}

Se recopilaron métricas de uso y se aplicaron mejoras basadas en el feedback recibido. A nivel técnico, 
se optimizó el rendimiento del servidor, se mejoró la interfaz de usuario y se resolvieron bugs 
importantes detectados por el equipo durante el desarrollo, así como otros que habían sido registrados 
previamente en el backlog. Durante esta etapa, el equipo destinó parte del tiempo al desarrollo del 
presente informe, cerrando así formalmente el ciclo del MVP.

\subsection{v0.8.0: No planificada}

Esta versión no forma parte del alcance definido del proyecto, pero se deja abierta como espacio para 
el crecimiento futuro del producto. El objetivo sería incorporar nuevas funcionalidades, aplicar mejoras 
basadas en el feedback de los usuarios y continuar con la optimización del sistema más allá del MVP.
